% ============================================================
% 分片：s07-classical  对应原文 PDF 第 87--96 页（印刷页 87--96）
% 负责章节：§7 The Classical Story + 7.1--7.4（§7 标题起）
% 译者：trans-s07   审校状态：初稿
% ============================================================

% ==== 原文 PDF 第 87 页（印刷页 87）====

\section{经典情形}

\begin{keytakeaway}
\textbf{核心物理图像：}本节把前面的几何图景翻译回经典（数论）语言：当曲线 $C$ 定义在有限域 $\mathbb{F}_q$ 上时，$C$ 上 ${}^{L}G$-局部系统的算术类比是 Weil 群同态 $\sigma:W_F\to{}^{L}G$，Hecke 本征层的类比则是 $G(\mathbb{A}_F)$ 的自守表示；几何内窥所对应的轨形点与分数膜，在经典一侧正是非平凡的 L-包与内窥群。
\textbf{数学工具：}有限域上曲线的 adele 环与局部域、Weil 群、Frobenius 共轭类与 Satake 对应、球 Hecke 代数、局部与整体 L-包、自同构群 $S_\sigma$ 及其分量群 $S_\sigma/S_\sigma^0$、迹公式（Labesse--Langlands）、双商与 $K$-不变向量空间。
\textbf{主要结论：}对 $SL_2$ 的非分歧情形，整体表示 $\bigotimes'_{x}\pi_x$ 自守当且仅当「取 $K''_x$ 的非分裂点」的个数 $\#S$ 为偶数，等价地，$S_\sigma/S_\sigma^0$ 在 $\prod_x S_{\sigma_x}/S_{\sigma_x}^0$ 的相应表示上平凡地作用；improper Hecke 算子恰在非平凡 L-包处消失，从而在函数论层面显现内窥。
\end{keytakeaway}

本节回顾自守形式的经典理论中内窥（endoscopy）与 L-包的框架，这一框架由 Labesse 与 Langlands [LL] 发现。我们将在下一节讨论上述几何图景对经典理论的潜在推论。

\subsection{局部与整体 Langlands 对应}

我们首先回顾经典与几何 Langlands 对应的一般框架。为简单起见，这里只限于非分歧情形。

几何 Langlands 对应（在一级近似下）预言：对 $C$ 上每个不可约 ${}^{L}G$-局部系统 $\mathcal{E}$，存在唯一的（在同构意义下）$\Bun_G$ 上以 $\mathcal{E}$ 为本征值的 Hecke 本征层 $(\mathcal{F}_{\mathcal{E}},(\alpha_V))$；这一概念已在 §6.3 详细讨论。

我们希望回顾几何 Langlands 对应与经典 Langlands 对应之间的关系，其中曲线 $C$ 定义在有限域 $k=\mathbb{F}_q$ 上。因此设 $C$ 是这样的曲线，$F$ 是 $C$ 的有理函数域。

例如，若 $C=\mathbb{P}^1$，则 $F$ 由分式 $f(z)/g(z)$ 组成，其中 $f(z)$ 与 $g(z)$ 是 $k$ 上的（关于变量 $z$ 的）多项式，没有公因子（且 $g(z)=z^m+\ldots$ 是首一的）。对 $C$ 的每个闭点 $x$，我们有局部域 $F_x\simeq k_x((t))$，其中 $t$ 是 $x$ 处的局部坐标，及其整数环 $\mathcal{O}_x\simeq k_x[[t]]$。这里 $k_x$ 是 $x$ 的剩余域，它一般是 $k$ 的有限扩张，因而同构于某个 $\mathbb{F}_{q^m}$（$m\geq 1$）。这个数称为 $x$ 的次数，记作 $\deg(x)$。例如，当 $C=\mathbb{P}^1$ 时，闭点对应于 $k[z]$ 中的不可约首一多项式，此外还有 $\infty$ 点。剩余域为 $k$ 的点对应于一次多项式 $z-a$，$a\in k$。一般的 $n$ 次不可约首一多项式 $P(z)$ 对应于 $C$ 的一个 $n$ 次闭点。域 $k_x$ 就是 $k[z]$ 对由 $P(z)$ 生成的主理想的商。

$F$ 的 adele 环 $\mathbb{A}_F$ 按定义是限制乘积
\[
\mathbb{A}_F = \prod_{x\in C}^{\prime} F_x,
\]
其中「restricted」一词（以及记号中的撇号）指的是：$\mathbb{A}_F$ 的元素是族 $(f_x)_{x\in C}$，其中对除有限多个 $x\in C$ 之外的所有 $x$ 都有 $f_x\in\mathcal{O}_x$。

设 $\mathrm{Gal}(\bar F/F)$ 为 $F$ 的 Galois 群，即 $F$ 的可分闭包 $\bar F$（由向 $F$ 添加所有系数在 $F$ 中的可分多项式的根得到）的、逐点固定 $F$ 的自同构群。我们有自然的同态 $\mathrm{Gal}(\bar F/F)\to\mathrm{Gal}(\bar k/k)$。群 $\mathrm{Gal}(\bar k/k)$ 由 Frobenius 自同构 $\mathrm{Fr}: y\mapsto y^q$ 拓扑生成，并同构于整数群 $\mathbb{Z}$ 的投射有限完备化 $\widehat{\mathbb{Z}}$。$\mathbb{Z}\subset\widehat{\mathbb{Z}}$ 在 $\mathrm{Gal}(\bar F/F)$ 中的原像就是 $F$ 的 Weil 群 $W_F$。

% ==== 原文 PDF 第 88 页（印刷页 88）====

Weil 群（更确切地说，其非分歧商）是 $C$ 的基本群的算术类比。因此，$C$ 上 ${}^{L}G$-局部系统的算术类比是同态\footnote{这里需要在域 $\overline{\mathbb{Q}}_\ell$ 上考虑 ${}^{L}G$，其中 $\ell$ 不整除 $q$，并考虑所谓的 $\ell$-进同态（例如见 [F1] 第 2.2 节）。}
\[
\sigma: W_F\to {}^{L}G.
\]

整体 Langlands 猜想粗略地预言：每个 $\sigma$ 对应 $G(\mathbb{A}_F)$ 的一个自守表示 $\pi(\sigma)$。这意味着它可以在商 $G(F)\backslash G(\mathbb{A}_F)$ 上的某个函数空间中被实现。\footnote{若 $\pi(\sigma)$ 不出现在「离散谱」中，这需要一些解释；但这里我们忽略这一技术问题。}

把这一整体 Langlands 猜想与局部猜想联系起来是有用的。回忆对 $C$ 的每个闭点 $x$（我们将写 $x\in C$），我们有局部域 $F_x\simeq k_x((t))$，它是 $F$ 的一个完备化。我们按与上面相同的方式定义它的 Weil 群 $W_{F_x}$，即作为 $\mathbb{Z}\subset\widehat{\mathbb{Z}}=\mathrm{Gal}(\bar k_x/k_x)$ 在同态 $\mathrm{Gal}(\bar F_x/F_x)\to\mathrm{Gal}(\bar k_x/k_x)$ 下于 $\mathrm{Gal}(\bar F_x/F_x)$ 中的原像。群 $W_{F_x}$ 可以实现为整体 Weil 群 $W_F$ 的子群，但不是典范地。然而，它在 $W_F$ 中的共轭类是典范的。因此，如上所述的同态 $\sigma: W_F\to {}^{L}G$ 的等价类给出同态的等价类
\[
\sigma_x: W_{F_x}\to {}^{L}G.
\]

局部 Langlands 猜想粗略地预言：对每个 $\sigma_x$，我们可以联系 $G(F_x)$ 的一个（光滑）不可约表示 $\pi_x$。取它们的限制张量积（见下文），我们得到 adele 群 $G(\mathbb{A}_F)$ 的不可约表示
\begin{equation}\tag{7.1}
\pi = \pi(\sigma) = \bigotimes_{x\in C}^{\prime}\pi_x
\end{equation}
整体与局部 Langlands 猜想之间的相容性，就是断言这个 $\pi$ 是自守的。此时 $\pi$ 就是对应于整体同态 $\sigma$ 的自守表示。概略地，
\[
\begin{array}{ccc}
\sigma_x & \xrightarrow{\ \text{local}\ } & \pi_x\\[3pt]
\sigma & \xrightarrow{\ \text{global}\ } & \pi = \bigotimes_{x\in C}^{\prime}\pi_x.
\end{array}
\]

这里假定同态 $\sigma$ 在 $C$ 上除有限多个点外非分歧。这意味着对除有限多个 $x\in C$ 之外的所有 $x$，同态 $\sigma_x: W_{F_x}\to{}^{L}G$ 经由商 $W_{F_x}\to\mathbb{Z}$ 分解。这个商 $\mathbb{Z}\subset\mathrm{Gal}(\bar k_x/k_x)$ 的生成元称为与 $x$ 相联系的 Frobenius 元素，记作 $\mathrm{Fr}_x$。由于 $\sigma_x$ 只在共轭意义下有定义，我们得到 $\sigma_x(\mathrm{Fr}_x)$ 给出 ${}^{L}G$ 中一个良定义的共轭类。人们相信，这样得到的共轭类总是半单的。于是，对除有限多个 $x\in C$ 之外的所有 $x$，我们得到 ${}^{L}G$ 中的一个半单共轭类 $\sigma_x(\mathrm{Fr}_x)$。

对应于非分歧 $\sigma_x$ 的不可约表示 $\pi_x$ 也是非分歧的，这意味着其不变向量子空间 $(\pi_x)^{G(\mathcal{O}_x)}\subset\pi_x$ 在极大紧子群 $G(\mathcal{O}_x)\subset G(F_x)$ 下是一维的。

% ==== 原文 PDF 第 89 页（印刷页 89）====

我们将对所有这样的 $x\in C$ 一次性固定一个非零的 $G(\mathcal{O}_x)$-不变向量 $v_x\in\pi_x$。那么公式 (7.1) 中的撇号（表示「限制张量积」）意味着：这个空间由形如 $\bigotimes_{x\in C}w_x$ 的向量张成，其中对除有限多个点 $x$ 之外的所有 $x$ 都有 $w_x=v_x$。

由于向量 $v_x$ 是 $G(\mathcal{O}_x)$-不变的，它是球 Hecke 代数（spherical Hecke algebra）的本征向量；该代数定义为 $G(F_x)$ 上紧支撑 $G(\mathcal{O}_x)$ 双不变函数构成的代数。由 Satake 对应，这个代数同构于表示环 $\mathrm{Rep}({}^{L}G)$（例如见 [F1] 第 5.2 节）。因此，每个 $V\in\mathrm{Rep}({}^{L}G)$ 对应于球 Hecke 代数的一个元素，我们记作 $T_{V,x}$。这些算子是 §6.3 所讨论的 Hecke 函子 $H_{V,x}$ 的函数论类比，作用在 $\pi_x$ 上。向量 $v_x\in\pi_x$ 是这一作用下的公共本征向量，本征值由 $\sigma_x(\mathrm{Fr}_x)$ 的共轭类记录。确切地说，我们有\footnote{这里及下文，我们在本式及其他类似公式的右端略去形如 $q_x^m$ 的因子。}
\begin{equation}\tag{7.2}
T_{V,x}\cdot v_x = \Tr(\sigma(\mathrm{Fr}_x),V)v_x.
\end{equation}
% [存疑: (7.2) 右端印作 $\sigma(\mathrm{Fr}_x)$；但本段文字称本征值由 $\sigma_x(\mathrm{Fr}_x)$ 的共轭类记录，且 (7.13) 用的是 $\sigma_x$。此处照录原文，按上下文应理解为同一共轭类的记号滥用。 来源：原书 p.89 英文原文；存疑待核（A）。]
% [Tutor Note: Satake 对应的定量形式。在 Satake 同构下 $T_{V,x}$ 对应于 $V$ 的矩阵系数沿 $\sigma_x(\mathrm{Fr}_x)$ 取值，故作用在球向量 $v_x$ 上给出本征值 $\Tr(\sigma_x(\mathrm{Fr}_x),V)$；换言之，半单共轭类 $\sigma_x(\mathrm{Fr}_x)$ 就是球表示 $\pi_x$ 的 Satake 参数。注意 (7.2) 印作 $\sigma$ 而非 $\sigma_x$，但本征值只依赖共轭类，故两种写法给出同一数。 来源：原书 p.89 上下文（A）。]

事实上，这一性质唯一确定了 $\pi_x$。
% [Tutor Note: 唯一性并不显然：它等价于非分歧局部 Langlands 对应的逆命题，即球表示由其在整个球 Hecke 代数上的作用（等价地，由所有 $V\in\mathrm{Rep}({}^{L}G)$ 的迹，亦即由 Satake 参数 $\sigma_x(\mathrm{Fr}_x)$ 的共轭类）唯一确定。作者在此后只用到这一形式事实：一旦局部因子 $\pi_x$ 被指定为某个 L-包的成员，$\pi_x$ 就被其本征值锁定。 来源：原书 p.89 上下文（A）。]

\subsection{L-包}

上面概述的局部与整体 Langlands 对应的图景对 $G=GL_n$ 是正确的（并且已被证明）。但对其他群需要作一些调整。最重要的调整是：一般而言，局部 Langlands 对应给每个 $\sigma_x$ 联系的不是单个不可约表示 $\pi_x$ 的等价类，而是一族 $\{\pi_{x,\alpha}\}_{\alpha\in A_x}$，称为（局部）L-包。通常，这对无穷多个点 $x\in C$ 都会发生；于是，在上述每个 L-包中各取一个表示，我们就得到 $G(\mathbb{A}_F)$ 的无穷多个表示。结果表明，一般而言它们并非都是自守的；而那些自守的表示，可以以不同的重数出现在 $G(F)\backslash G(\mathbb{A}_F)$ 上的函数空间中。对那些经由与 $G$ 的内窥群对偶的群 ${}^{L}H$ 分解的同态 $\sigma: W_F\to{}^{L}G$，这个重数有一个漂亮的组合公式（见 [LL, Ko]）。

这一现象最初是 Labesse 与 Langlands 在 [LL] 中对 $G=SL_2$ 的情形发现的。下面简要总结他们在非分歧情形下的结果。

在此情形 ${}^{L}G=SO_3=PGL_2$。设 $E$ 是 $F$ 的非分歧二次扩张。这是我们的曲线 $C$ 的一个二次非分歧覆盖 $C'$ 的函数域 $k(C')$。设 $\mu$ 是一个特征标，即 Weil 群 $W_E$ 的一维（$\ell$-进）表示。群 $W_F$ 以 $W_E$ 为正规子群，且商同构于 $\mathrm{Gal}(E/F)=\{1,\tau\}$。设
\[
\sigma: W_F\to PGL_2
\]
是由 $\mu$ 诱导的 $W_F$ 的二维表示 $\tilde\sigma$ 的射影化，
\[
\tilde\sigma = \operatorname{Ind}_{W_E}^{W_F}\mu.
\]
显然 $\sigma$ 的像包含在子群 $O_2\subset SO_3=PGL_2$ 中。下文我们假定像不属于子群 $\mathbb{Z}_2\times\mathbb{Z}_2\subset O_2$

% ==== 原文 PDF 第 90 页（印刷页 90）====

也不属于 $O_2$ 的连通分支 $SO_2$。在此情形，$\sigma$ 的像的中心化子等于 $\mathbb{Z}_2$。相应的同态 $\sigma: W_F\to PGL_2$ 就是我们在几何设定中于上文考虑过的 $SO_3=PGL_2$-局部系统的算术类比。

下面考察对应于 $\sigma$ 的局部同态
\[
\sigma_x: W_{F_x}\to PGL_2.
\]
有两种可能：
\begin{itemize}[leftmargin=2em,itemsep=2pt,topsep=2pt]
\item 点 $x$ 在 $E/F$ 中分裂。这意味着 $x$ 在 $C'$ 中的原像由两个闭点组成，记作 $y_1(x)$ 与 $y_2(x)$，且它们的剩余域都等于 $x$ 的剩余域。在此情形，$W_{F_x}$ 同构于 $W_{E_{y_1(x)}}$ 与 $W_{E_{y_2(x)}}$ 两者；当它们实现为 $W_F$ 的子群时，彼此共轭：
\[
\tilde\tau W_{E_{y_1(x)}}\tilde\tau^{-1} = W_{E_{y_2(x)}},
\]
其中 $\tilde\tau$ 在同态 $W_F\to\mathrm{Gal}(E/F)$ 下投影到非平凡元素 $\tau\in\mathrm{Gal}(E/F)$。
\item 点 $x$ 非分裂。则 $x$ 在 $C'$ 中的原像由一个闭点组成，记作 $y(x)$；其剩余域 $k_{y(x)}$ 是 $x$ 的剩余域 $k_x$ 的二次扩张。于是，若 $k_x\simeq\mathbb{F}_{q^n}$，则 $k_{y(x)}\simeq\mathbb{F}_{q^{2n}}$。在此情形，$W_{F_x}$ 以 $W_{E_{y(x)}}$ 为正规子群，且商同构于 $\mathrm{Gal}(E_{y(x)}/F_x)\simeq\mathbb{Z}_2$。
\end{itemize}

在第一种情形，$\sigma_x$ 等价于 $W_{F_x}$ 的二维表示 $\tilde\sigma_x$ 的射影化，后者由下述公式定义（这里把 $W_{F_x}$ 与 $W_{E_{y_1(x)}}$ 等同）：
\begin{equation}\tag{7.3}
g\mapsto\begin{pmatrix}\mu_{y_1(x)}(g)&0\\ 0&\mu_{y_2(x)}(\tilde\tau g\tilde\tau^{-1})\end{pmatrix},
\qquad g\in W_{E_{y_1(x)}},
\end{equation}
其中 $\mu_{y_i(x)}$ 表示 $\mu$ 在 $W_{E_{y_i(x)}}$ 上的限制。

在第二种情形，$\sigma_x$ 同构于二维诱导表示的射影化
\begin{equation}\tag{7.4}
\tilde\sigma_x = \operatorname{Ind}_{W_{E_{y(x)}}}^{W_{F_x}}\mu_{y(x)}.
\end{equation}
% [Tutor Note: 诱导表示的显式矩阵。在 (7.3) 中把 $W_{F_x}$ 与 $W_{E_{y_1(x)}}$ 等同后，对 $g\in W_{E_{y_1(x)}}$ 矩阵是对角的（两个特征标 $\mu_{y_1(x)}$ 与 $\mu_{y_2(x)}$），而对换位陪集中的元素则给出反角矩阵，这正是二维诱导表示的标准矩阵形式；因此在分裂点处 $\sigma_x$ 的像落在 $O_2$ 型子群中。$\tilde\tau$ 的选取相差一个 $W_{E_{y_1(x)}}$ 元素，换 $\tilde\tau$ 相当于换基，矩阵只相差共轭，故 (7.3) 只在射影化意义下定义 $\sigma_x$。 来源：原书 p.90 上下文（A）。]

正如 Labesse--Langlands 所证明的，局部 L-包包含 $SL_2(F_x)$ 的一个或两个不可约表示。在这两种情形，它们都作为 $GL_2(F_x)$ 的同一个不可约表示的不可约直和项出现。

从今以后，我们专注于从 Weil 群 $W_F$ 到 $PGL_2$ 的非分歧同态。我们已经假定覆盖 $C'\to C$ 非分歧；从今以后还假定 $W_E$ 的特征标 $\mu$ 也非分歧。那么我们的 $\sigma$ 非分歧，因而各局部同态 $\sigma_x: W_{F_x}\to PGL_2$ 也非分歧。\footnote{然而要注意，存在分歧的特征标 $\mu$，使得相应的同态 $\sigma$ 非分歧。要求 $\mu$ 非分歧，类似于在几何理论中考虑 $w_2=0$ 的 $PGL_2$-局部系统。}

% ==== 原文 PDF 第 91 页（印刷页 91）====

因此 $\sigma_x$ 由 Frobenius 共轭类 $\mathrm{Fr}_x$ 在 $PGL_2$ 中的像 $\sigma_x(\mathrm{Fr}_x)$ 决定。有两种可能：若 $\sigma_x(\mathrm{Fr}_x)$ 是元素
\begin{equation}\tag{7.5}
\begin{pmatrix}0&1\\ 1&0\end{pmatrix}\in PGL_2,
\end{equation}
的共轭类，该元素的中心化子有两个连通分支，则 L-包由 $SL_2(F_x)$ 的两个不可约表示组成；否则中心化子连通，L-包由单个不可约表示组成。

前一情形可以以两种方式出现。首先，设点 $x$ 非分裂。在诱导表示 (7.4) 中取基 $\{1,\mathrm{Fr}_x\}$。则我们有
\[
\tilde\sigma_x(\mathrm{Fr}_x)=\begin{pmatrix}0&\mu(\mathrm{Fr}_{y(x)})\\ 1&0\end{pmatrix}.
\]
因此我们得到 $\sigma_x(\mathrm{Fr}_x)$ 是 (7.5) 的共轭类。第二种可能是 $x$ 分裂，且 $\mu$ 在 $\mathrm{Fr}_{y_1(x)}$ 与 $\mathrm{Fr}_{y_2(x)}$ 上的值相差一个负号。在此情形，$\tilde\sigma_x(\mathrm{Fr}_x)$ 共轭于
\begin{equation}\tag{7.6}
\mu(\mathrm{Fr}_{y_i(x)})\begin{pmatrix}1&0\\ 0&-1\end{pmatrix}\sim\begin{pmatrix}1&0\\ 0&-1\end{pmatrix}\in PGL_2,\qquad i=1,2,
\end{equation}
因此 $\sigma_x(\mathrm{Fr}_x)$ 在 $PGL_2$ 中仍是 (7.5) 的共轭类。然而下面将看到，第二种情形在内窥中不起重要作用。
% [Tutor Note: 易错预警。L-包的大小由 $\sigma_x(\mathrm{Fr}_x)$ 在 $\mathrm{PGL}_2$ 中的中心化子的连通分支数决定：(7.5) 是行列式为 $-1$ 的反射，它在 $\mathrm{PGL}_2$ 中的中心化子（无标量歧义地）就是 $O_2$，有两个连通分支，故 L-包有两个成员；中心化子连通时只有一个成员。注意分裂点（本征值之比为 $-1$）与非分裂点都给出 (7.5) 的共轭类，但整体元素 (7.10) 落在这类局部中心化子群的不同连通分支中（分裂点落在单位分支，非分裂点落在另一分支），这正是 (7.9) 在两种情形下平凡与否的根源，也是内窥判据的几何来源。 来源：原书 p.91 上下文（A）。]

在这两种情形，相应的 L-包都由 $SL_2(F_x)$ 的两个不可约表示 $\pi'_x$ 与 $\pi''_x$ 组成。我们将在下一节更详细地讨论它们。

于是我们看到，在 $C$ 的无穷多个点上，对于张量积 (7.1) 中 $SL_2(F_x)$ 的两个可能不可约表示 $\pi_x$，我们有一个二选一的选择。然而结果表明，这些选择中大约只有一半——即在 $C$ 的非分裂点处（相对于覆盖 $C'\to C$）的选择——给出 $SL_2(\mathbb{A}_F)$ 的自守表示。更确切地说，如果我们在 $C$ 的所有非分裂点（除一个之外，设为 $y$）处选取局部 L-包中的特定表示 $\pi_x$，那么在剩余点 $y$ 处，局部 L-包的两个表示 $\{\pi'_y,\pi''_y\}$ 中只有一个能使张量积 $\bigotimes_{x\neq y}^{\prime}\pi_x$ 成为 $SL_2(\mathbb{A}_F)$ 的自守表示 $\pi$。

关于哪些不可约表示的集合其张量积是自守的，确切描述将在下一节给出。这一描述最初由 Labesse 与 Langlands [LL] 利用迹公式（trace formula）给出。在附录中，我们将利用 Weil [W2] 与 Jacquet--Langlands [JL] 关于 Whittaker 函数 Fourier 变换的 Hecke 本征函数显式公式，给出这一描述的另一种推导。

\subsection{不变向量空间}

下文的目标是从经典设定过渡到几何设定。这一方向的第一步，是把 $G(\mathbb{A}_F)$ 的无穷维表示 $\pi=\bigotimes_{x\in C}^{\prime}\pi_x$ 缩减为有限维向量空间 $\pi^K$，即 $K$-不变向量构成的空间，其中 $K$ 是 $G(\mathbb{A}_F)$ 的紧子群。我们取 $K$ 为乘积
\[
K=\prod_{x\in C}K_x
\]

% ==== 原文 PDF 第 92 页（印刷页 92）====

其中 $K_x\subset G(F_x)\simeq G((t))$ 是紧子群。一个典型的例子是子群 $G(\mathcal{O}_x)=G[[t]]$。正如 §7.1 所讨论的，$\pi$ 中任何向量都在下述子群下不变：该子群是除有限多个 $x$ 之外所有 $G(\mathcal{O}_x)$ 的乘积。若 $\pi$ 自守，则 $\pi^K$ 可实现为双商 $G(F)\backslash G(\mathbb{A}_F)/K$ 上的函数空间中的空间，这些函数对所有满足 $K_x=G(\mathcal{O}_x)$ 的 $x\in C$ 都是 Hecke 本征函数。这个双商有一个几何解释：它是 $C$ 上带抛物（或级）结构的 $G$-丛模叠的 $\mathbb{F}_q$-点集，这些结构位于 $K_x$ 非极大紧的那些点 $x\in C$ 处（例如，取 $K_x$ 为 Iwahori 子群对应于在 $G$-丛于 $x$ 处的纤维中固定一个 Borel 约化，等等）；这一事实将在下文用于联系经典与几何 Langlands 猜想。
% [Tutor Note: 物理图像。双商 $G(F)\backslash G(\mathbb{A}_F)/\prod_x K_x$ 的 $\mathbb{F}_q$-点正是 $C$ 上带相应级结构的 $G$-丛模叠的 $\mathbb{F}_q$-点：在 $K_x=G(\mathcal{O}_x)$ 处丛无额外结构，在 $K_x$ 为 Iwahori 子群处则固定该点纤维中的一条 Borel 约化（等价于一个完整旗）。这正是 §8 把 Hecke 本征函数翻译为 $\Bun_G$ 上 Hecke 本征层的入口：Hecke 算子的双陪集在 $C$ 上对应「添加一个修正点」，其几何版本就是 Hecke 函子。 来源：原书 p.92 上下文（A）。]

在非分歧情形，非平凡 L-包的存在与下述事实有关：$G(F_x)$ 的极大紧子群有若干不等价的选取。例如，对 $G=SL_2$ 有两种不等价的选取：其一是 $SL_2[[t]]$，其二是
\begin{equation}\tag{7.7}
\begin{pmatrix}t&0\\ 0&1\end{pmatrix}SL_2[[t]]\begin{pmatrix}t^{-1}&0\\ 0&1\end{pmatrix}
\end{equation}
（它们在 $GL_2((t))$ 中共轭，但在 $SL_2((t))$ 中不共轭）。我们将这两个子群分别记作 $K'$ 与 $K''$。
% [Tutor Note: 易错预警。(7.7) 的两个矩阵因子分别是 $g$ 与 $g^{-1}$（$g=\mathrm{diag}(t,1)$），故 $K''=gK'g^{-1}$；二者在 $GL_2((t))$ 中共轭而在 $SL_2((t))$ 中不共轭。证明：若 $h\in SL_2((t))$ 把 $K'$ 共轭到 $K''$，则 $g^{-1}h$ 属于 $K'$ 在 $SL_2((t))$ 中的正规化子，而后者就是 $K'$，于是 $\det(g^{-1}h)\in 1+tk[[t]]$，但 $\det(g^{-1}h)=t^{-1}$，矛盾。注意与 (7.14) 的积分区域 $M_x=gK'g$（两个因子都是 $g$）不同：这里定义的是子群，那里定义的是积分区域，切勿混淆。 来源：原书 p.92 上下文（A）。]

朴素地看，似乎 $K'$ 是更可取的选择，但那是因为我们默许把曲线 $C$ 上的「常值」群概形 $SL_2$ 取作初始数据。然而，我们的初始数据实际上应当是域 $F$（$C$ 上有理函数域）上的群 $SL_2$，换言之，是 $C$ 的一般点上的群概形。把它延拓为整个 $C$ 上的群概形有很多方式，并且我们有这样的自由：延拓方式使得在 $x\in C$ 处，作为圆盘 $D_x$ 上截面群得到的 $SL_2(F_x)=SL_2((t))$ 的子群共轭于 $K'$ 或 $K''$。\footnote{作为类比，考虑 $C\backslash\{x_1,\ldots,x_m\}$ 上线丛的数据——它可以以许多不同方式延拓为整条曲线 $C$ 上的线丛。}这就清楚表明，在 $K'$ 与 $K''$ 之间先验地没有典范的选取。

然而，下文中我们将固定群概形 $SL_2$ 从 $C$ 的一般点到整个 $C$ 的一个特定延拓，即「常值」群概形 $C\times SL_2$。于是，在每个点 $x\in C$ 处我们有一个优选的紧子群 $K'_x=SL_2[[t]]$。

$SL_2(F_x)$ 的极大紧子群选取上的这种不明确性，与下面的事实密切相关：非分歧局部 L-包 $\{\pi'_x,\pi''_x\}$ 的结构。确切地说，在适当记号下，不变向量空间 $(\pi'_x)^{K'_x}$ 与 $(\pi''_x)^{K''_x}$ 是一维的，而 $(\pi'_x)^{K''_x}=(\pi''_x)^{K'_x}=0$。

让我们回到前一节的设定。于是，给定 $C$ 的非分歧二次覆盖 $C'$ 以及 $W_E$ 的非分歧特征标 $\mu$，其中 $E$ 是 $C'$ 上的函数域。设 $\sigma$ 是按上述方式定义的同态 $W_F\to PGL_2$。现在，在 $\sigma_x(\mathrm{Fr}_x)$ 共轭于 (7.5) 的每个点 $x$ 处，我们固定两个极大紧子群 $K_x$ 之一。我们假定在除有限多个点之外的所有点处 $K_x=K'_x=SL_2[[t]]$。

于是，我们有一个 $SL_2(F_x)$ 的不可约表示 $\pi_x$，使得在所有这些点处 $(\pi_x)^{K_x}$ 是一维的。用 $S$ 表示 $C$ 的（相对于覆盖 $C'\to C$）非分裂点 $x\in C$ 中满足 $K_x=K''_x$ 的有限集。设 $\#S$ 为其基数。下述陈述属于 [LL]（另一种证明见附录）。

% ==== 原文 PDF 第 93 页（印刷页 93）====

\begin{namedthm}{7.1}
表示
\begin{equation}\tag{7.8}
\bigotimes_{x\in C}^{\prime}\pi_x
\end{equation}
是 $SL_2(\mathbb{A}_F)$ 的自守表示，当且仅当 $\#S$ 为偶数。
\end{namedthm}

Labesse 与 Langlands 以下述简洁形式表述这一条件，它允许推广到其他群（见 [Ko]）。用 $S_\sigma$ 表示同态 $\sigma: W_F\to PGL_2$ 的自同构群，即 $\sigma$ 的像在 $PGL_2$ 中的中心化子。\footnote{这是上一节讨论的群 $\Gamma$ 的算术类比。为了遵循标准记号，我们把它记作 $S_\sigma$。}设 $S_\sigma^0$ 为其连通分支。同样，对每个 $x\in C$，设 $S_{\sigma_x}$ 是 $\sigma_x: W_F\to PGL_2$ 的自同构群，$S_{\sigma_x}^0$ 为其连通分支。我们有自然同态 $S_\sigma\to S_{\sigma_x}$ 与 $S_\sigma^0\to S_{\sigma_x}^0$，从而有同态
\begin{equation}\tag{7.9}
S_\sigma/S_\sigma^0\to S_{\sigma_x}/S_{\sigma_x}^0.
\end{equation}
% [原文疑误: 此处印作 $\sigma_x: W_F\to PGL_2$，但按定义 $\sigma_x$ 定义在局部 Weil 群上（比较本页上一段及 (7.13) 的用法）；其定义域应为 $W_{F_x}$。译文照录原文，未作改动。 来源：原书 p.93 英文原文与译本对照（A）。]

在我们的情形，对这里所考虑的类型中的一般 $\sigma$，我们有 $S_\sigma=S_\sigma/S_\sigma^0=\mathbb{Z}_2$，它（相对于诱导表示中的自然基）由元素
\begin{equation}\tag{7.10}
\begin{pmatrix}1&0\\ 0&-1\end{pmatrix}
\end{equation}
生成，该元素属于 $PGL_2$（它生成 $O_2\subset PGL_2$ 的中心化子）。

现在考虑局部群。若 $x$ 是 $C$ 的分裂点，且 $\sigma_x(\mathrm{Fr}_x)$ 的本征值之比不等于 $-1$，则 $S_{\sigma_x}$ 是连通的环面，故 $S_{\sigma_x}/S_{\sigma_x}^0$ 平凡。接下来考虑本征值之比等于 $-1$ 的分裂点情形。此时群 $S_{\sigma_x}$ 是 $PGL_2$ 的子群 $O_2$，即 (7.10) 的中心化子（见公式 (7.6)），故 $S_{\sigma_x}/S_{\sigma_x}^0=\mathbb{Z}_2$。但整体自同构 (7.10) 落在 $S_{\sigma_x}$ 的连通分支 $S_{\sigma_x}^0$ 中，因此在此时同态 (7.9) 平凡。最后，若 $x$ 非分裂，则 $S_{\sigma_x}$ 是 (7.5) 的中心化子，它也同构于 $O_2$。然而，此时元素 (7.10) 落在另一个连通分支中，因此在此时同态 (7.9) 非平凡。
% [原文疑误: 本句讨论的是局部群 $S_{\sigma_x}$，且下文的结论依赖 $S_{\sigma_x}/S_{\sigma_x}^0\simeq\mathbb{Z}_2$；此处印作 $S_\sigma/S_\sigma^0=\mathbb{Z}_2$，应为 $S_{\sigma_x}/S_{\sigma_x}^0=\mathbb{Z}_2$（前一句的主语即 $S_{\sigma_x}$，且紧接着论及 $S_{\sigma_x}^0$）。 来源：原书 p.93 英文原文与译本对照（A）。]

[LL] 的想法是：来自局部 L-包的 $SL_2(F_x)$ 的不可约表示应当由 $S_{\sigma_x}/S_{\sigma_x}^0$ 的不可约表示来标记。于是，若这个群平凡，则 L-包中只有一个不可约表示。若这个群同构于 $\mathbb{Z}_2$，则有两个。按我们的约定，具有非平凡 $K'_x$-不变向量空间的表示将对应于 $\mathbb{Z}_2$ 的平凡表示，而具有非平凡 $K''_x$-不变向量的表示将对应于 $\mathbb{Z}_2$ 的符号表示。现在，张量积 (7.8) 给出群 $\prod_{x\in C}S_{\sigma_x}/S_{\sigma_x}^0$ 的一个不可约表示，其中除有限多个因子外所有因子都平凡地作用。

% ==== 原文 PDF 第 94 页（印刷页 94）====

于是定理 7.1 可以重新表述为：(7.8) 是自守的，当且仅当 $S_\sigma/S_\sigma^0$ 经由对角同态
\[
S_\sigma/S_\sigma^0\to\prod_{x\in C}S_{\sigma_x}/S_{\sigma_x}^0
\]
平凡地作用在 $\prod_{x\in C}S_{\sigma_x}/S_{\sigma_x}^0$ 的相应表示上。
% [Tutor Note: 推导补足，说明定理 7.1 的群论重述如何给出奇偶条件。每个非分裂点 $x$ 处，L-包的两个成员由 $S_{\sigma_x}/S_{\sigma_x}^0\simeq\mathbb{Z}_2$ 的两个一维表示（平凡表示与符号表示）标记，按约定取 $K'_x$ 对应平凡表示、取 $K''_x$ 对应符号表示；分裂点处 (7.9) 的像平凡，故局部选择无关紧要。于是 $S_\sigma/S_\sigma^0=\mathbb{Z}_2$ 在对角嵌入下的像在第 $x$ 个分量上等于符号表示当且仅当 $x\in S$；该整体元素作用平凡当且仅当符号出现的次数 $\#S$ 为偶数，这正是定理 7.1。 来源：原书 p.94 上下文（A）。]

注意，按上面的讨论，若 $x$ 分裂，则同态 (7.9) 的像平凡，即使群 $S_{\sigma_x}/S_{\sigma_x}^0$ 非平凡。因此，我们可以取与该点相联系的局部 L-包中 $SL_2(F_x)$ 的两个不可约表示中的任一个作为 $\pi_x$，并且在两种情形下相应的表示 (7.8) 将同时自守或同时不自守。在这个意义上，分裂点不影响表示 (7.8) 的自守性；与之不同，对非分裂点而言，我们取 L-包两个成员中的哪一个作为 (7.8) 的局部因子是至关重要的。

现在假设 $\#S$ 为偶数，因而表示 (7.8) 是自守的。则一维向量空间
\begin{equation}\tag{7.11}
\bigotimes_{x\in C}(\pi_x)^{K_x}
\end{equation}
可以在
\begin{equation}\tag{7.12}
SL_2(F)\backslash SL_2(\mathbb{A}_F)\Big/\prod_{x\in C}K_x
\end{equation}
上的函数空间中被实现。此外，(7.11) 中任何非零向量都给出 (7.12) 上的一个 Hecke 本征函数 $f$，其本征值由共轭类 $\sigma_x(\mathrm{Fr}_x)$ 指定。这意味着它是与 $PGL_2$ 的伴随表示 $W$ 及点 $x\in C$ 相对应的 Hecke 算子 $T_{W,x}$ 的本征函数，即
\begin{equation}\tag{7.13}
T_{W,x}\cdot f = \Tr(\sigma_x(\mathrm{Fr}_x),W)f,
\end{equation}
其中 $\mathrm{Fr}_x$ 是 $W_F$ 中对应于 $x$ 的 Frobenius 共轭类。这里 $T_{W,x}$ 是 $SL_2(F_x)$ 上紧支撑 $K_x$ 双不变函数的球 Hecke 代数的一个生成元。对 $K_x$ 的两种选取，这个代数都典范同构于 $\mathrm{Rep}(PGL_2)$；在此同构下，$T_{W,x}$ 对应于 $PGL_2$ 的伴随表示的类（关于 Hecke 性质的更多细节见下文 §8.1）。

最后，假设 $\#S$ 为奇数。则表示 (7.8) 不是自守的。因此一维向量空间 (7.11) 不可能被实现为 (7.12) 上的函数空间中的空间。换言之，(7.12) 上任何满足 (7.13) 的函数都必然恒等于零。

这应与一般情形形成对照：此时 $\sigma: W_F\to PGL_2$ 的像具有平凡的中心化子（回忆我们这里只考虑非分歧同态 $\sigma$）。在此情形，群 $S_\sigma/S_\sigma^0$ 平凡，所以所有表示 (7.8)（其中局部因子 $\pi_x$ 是局部 L-包中对应于 $\sigma_x$ 的任意表示）都将自守。若 $\sigma_x(\mathrm{Fr}_x)$ 是一般的，则 $S_{\sigma_x}/S_{\sigma_x}^0$ 平凡，且相应的 L-包包含 $SL_2(F_x)$ 的一个不可约表示 $\pi_x$。这个 $\pi_x$ 在 $K'_x$ 与 $K''_x$ 下的不变向量空间都是一维的。也可能对某个 $x\in C$，$\sigma_x(\mathrm{Fr}_x)$ 处于 (7.8) 的共轭类中。此时 $S_{\sigma_x}/S_{\sigma_x}^0$
% [原文疑误: 此处印作 (7.8)，但 (7.8) 是一个表示（张量积），没有「共轭类」可言；结合上文（(7.5) 之后两段）此处显然应为 (7.5)，即 $\sigma_x(\mathrm{Fr}_x)$ 属于矩阵 $\begin{pmatrix}0&1\\1&0\end{pmatrix}$ 的共轭类。译文照录原文。 来源：原书 p.94 英文原文与译本对照（A）。]

% ==== 原文 PDF 第 95 页（印刷页 95）====

等于 $\mathbb{Z}_2$，且 L-包包含两个不可约表示，其中一个在 $K'_x$ 下具有非零不变向量，另一个在 $K''_x$ 下具有非零不变向量。然而，把其中任一个作为 $\pi$ 在 $x$ 处的局部因子代入，我们都会得到 $SL_2(\mathbb{A}_F)$ 的自守表示。

于是我们发现，对一般的 $\sigma$，任意选取 $K_x$（即 $K_x=K'_x$ 或 $K_x=K''_x$），相应的 (7.12) 上满足 (7.12) 的 Hecke 本征函数空间都是一维的。
% [原文疑误: 句首的第一个 (7.12) 应为 (7.13)，即 Hecke 本征函数性质；本段讨论的是满足 (7.13) 的 Hecke 本征函数在双商 (7.12) 上构成的空间，两次出现 (7.12) 中前者显然是笔误。译文照录原文。 来源：原书 p.95 英文原文与译本对照（A）。]

\subsection{improper Hecke 算子}

在几何理论中，除了对应于 $SO_3=PGL_2$ 的三维伴随表示的「proper」Hecke 算子之外，考虑对应于该群的二维射影表示的「improper」算子是很有用的。这些算子在经典理论中也有对应物。它们的定义如下。给定点 $x\in C$ 及 $x$ 处的坐标 $t$，算子 $\widetilde{T}_x$ 是从双商 (7.12)（其中 $K_x=K'_x=SL_2[[t]]$）上的函数到同一个双商（但把 $K'_x$ 换成式 (7.7) 中的子群 $K''_x$）上的函数的积分算子。它由公式
\begin{equation}\tag{7.14}
(\widetilde{T}_x\cdot f)(g)=\int_{M_x}f(gh)\,dh,
\end{equation}
给出，其中
\[
M_x=\begin{pmatrix}t&0\\ 0&1\end{pmatrix}SL_2[[t]]\begin{pmatrix}t&0\\ 0&1\end{pmatrix}.
\]
% [Tutor Note: 推导补足，验证 (7.14) 确实把 $K'_x$-右不变函数送到 $K''_x$-右不变函数。写 $k''=gk'g^{-1}\in K''$，则 $k''M_x=gk'g^{-1}\cdot gK'g=gk'K'g=gK'g=M_x$；于是在积分 (7.14) 中作变量替换 $h\mapsto k''h$ 不改变积分区域，从而 $(\widetilde{T}_x f)(gk'')=(\widetilde{T}_x f)(g)$。注意积分区域必须写成 $gK'g$（两个因子都是 $g$），若写成 $gK'g^{-1}=K''$ 就不具有这一性质。 来源：原书 p.95 上下文（A）。]

现在假设我们有一个 Hecke 本征函数 $f'$（相应地 $f''$）在 (7.12) 上，其中 $K_x=K'_x$（相应地 $K_x=K''_x$），且 $K_y$（$y\neq x$）相同，并满足（proper）Hecke 本征函数性质 (7.13)。这两个函数各自在相差一个标量意义下唯一确定。假设我们可以把 $\sigma$ 提升为非分歧同态 $\tilde\sigma: W_F\to GL_2$。那么我们可以把函数 $f'$ 与 $f''$ 规范化，使得二者都等于对应于 $\tilde\sigma$ 的 $GL_2(\mathbb{A}_F)$ 的 Hecke 本征函数 $\tilde f$ 在相应双商 (7.12) 上的限制（更多细节见附录）。由于 $GL_2$ 的所有 Hecke 算子彼此交换，我们发现函数 $\widetilde{T}_x\cdot f'$ 是 (7.12)（其中 $K_x=K''_x$）上的函数，并且也满足 (7.13)。此外，我们有
\begin{equation}\tag{7.15}
\widetilde{T}_x\cdot f' = \Tr(\tilde\sigma_x(\mathrm{Fr}_x),V)f'',
\end{equation}
其中 $V$ 是 $GL_2$ 的二维表示。于是，$\widetilde{T}_x$ 是 $SL_2(\mathbb{A}_F)$ 的两个 Hecke 本征函数空间之间的交织算子，当且仅当迹 $\Tr(\tilde\sigma_x(\mathrm{Fr}_x),V)$ 非零。
% [Tutor Note: 物理图像。$\Tr(\tilde\sigma_x(\mathrm{Fr}_x),V)$ 是二维表示 $V$ 在 $\tilde\sigma_x(\mathrm{Fr}_x)$ 上的特征标。当 $\sigma_x(\mathrm{Fr}_x)$ 为 (7.5) 的共轭类时，其在 $\mathrm{PGL}_2$ 中的代表元本征值为 $1$ 与 $-1$，故任何提升 $\tilde\sigma_x(\mathrm{Fr}_x)$ 的两本征值之和为 $0$，迹因子消失，(7.15) 给出 $\widetilde{T}_x\cdot f'=0$。这正是「improper Hecke 算子在非平凡 L-包处为零」的机制，与 §8.7 中几何上分数膜只在非平凡内窥局部系统处出现的现象相呼应。 来源：原书 p.95 上下文（A）。]

但这个迹为零，恰好是在 $\sigma_x(\mathrm{Fr}_x)$ 为 (7.5) 的共轭类时；而这正是我们在 $x$ 处有非平凡 L-包的特殊情形！在此情形，$f'$ 与 $f''$ 对应于 $SL_2(F_x)$ 的两个不同构的表示（其中之一可能自守，另一个则不然）。公式 (7.15) 表明，在此情形 $\widetilde{T}_x\cdot f'=0$。于是，improper Hecke 算子给了我们观察 L-包非平凡性的另一种方式。它们突出了在 $K_x=K'_x$

% ==== 原文 PDF 第 96 页（印刷页 96）====

与 $K_x=K''_x$ 两种情形下，双商 (7.12) 上的两个 Hecke 本征函数空间之间的差异：在内窥的 $\sigma$ 情形，这两个空间之一可能是一维的，而另一个可能等于零。这一现象的类比或许可以在几何上观察到，如下文 §8.7 所示。
